\clearpage
\begin{trapbox}
    \textbf{\textcolor{CTrap}{易错点一：只证明“某个数取不到”，就直接写完整值域}}

    \textbf{错误想法：} 既然 $y\ne a/c$，其他数一定都能取到。

    \textbf{正确做法：} 还要说明“任给 $t\ne a/c$，都能找到合法的 $x$”。本结论已用 $x=(b-dt)/(ct-a)$ 证明；若题目另外限制 $x$，就必须重新分析，不能照搬。

    \medskip
    \begin{minipage}{\linewidth}
\textbf{\textcolor{CTrap}{易错点二：忘记检查 $ad\ne bc$}}
\begin{center}\begin{tikzpicture}[font=\small,>=stealth,line join=round]
\draw[gray!45,rounded corners=3pt] (0,0) rectangle (14.2,9.4);
\node at (7.1,8.8) {行列式为零时，图像退化为带缺口的水平线};

\begin{scope}[shift={(4.3,2.8)},scale=1.2]
\draw[->] (-2.9,0)--(2.8,0) node[right]{$x$};\draw[->] (0,-.4)--(0,3.2) node[above]{$y$};\draw[blue,thick] (-2.8,2)--(2.6,2);\draw[dashed] (-2,0)--(-2,2);\node[below] at (-2,0){$-2$};\node[above left] at (0,2){$2$};
\draw[blue,fill=white,thick] (-2,2) circle(2.8pt);
\end{scope}
\node[align=left,anchor=west] at (8.4,5){$\dfrac{2x+4}{x+2}=2$\\[10pt] 定义域仍排除 $-2$\\[10pt] 值域只有 $\{2\}$};
\node[text width=13.2cm,align=center] at (7.1,.65) {一个输入点被排除，不会让常数输出消失。};
\end{tikzpicture}
\end{center}
\end{minipage}



    \textbf{反例：} $\dfrac{2x+4}{x+2}=2$（$x\ne-2$）。这时分子就是分母的 $2$ 倍，值域是 $\{2\}$，\textbf{不是} $y\ne2$。

    \textbf{记法：} 分离常数后，如果“剩下的数”恰好为 $0$，答案立即变成常数值域。

    \medskip
    \textbf{\textcolor{CTrap}{易错点三：把分子拆错，或写错符号}}

    \textbf{常见错误：} $\dfrac{ax+b}{cx+d}=\dfrac ac+\dfrac{b}{cx+d}$。

    \textbf{正确写法：}
    \[
        ax+b=\frac ac(cx+d)+\frac{bc-ad}{c},\qquad
        y=\frac ac+\frac{bc-ad}{c(cx+d)}.
    \]
    \textbf{自检：} 把右边通分，还原不出原分式，就说明拆错了。

    \medskip
    \textbf{\textcolor{CTrap}{易错点四：把函数的定义域和完整值域混淆}}

    \textbf{记住两件不同的事：} $x=-d/c$ 让分母为 $0$，所以不能作为输入；$y=a/c$ 是这类函数无法得到的输出。两者的符号和位置不能颠倒。

    \medskip
    \textbf{\textcolor{CTrap}{易错点五：题目只允许一部分 $x$，仍直接套公式}}

    \textbf{反例：} $y=2+\dfrac3{x-1}$，若 $x\ge2$，则 $0<3/(x-1)\le3$，因此 $2<y\le5$，并非除 $2$ 外的全部实数。

    \medskip
    \textbf{\textcolor{CTrap}{易错点六：把 $c=0$ 的情况也当作本结论}}

    \textbf{注意：} 如果 $c=0$、$d\ne0$，原式变成 $y=(ax+b)/d$，是一次函数（$a\ne0$）或常数函数（$a=0$），不能使用 $a/c$；如果 $c=d=0$，分母一直为 $0$，根本没有合法的 $x$。
\end{trapbox}
